%%% Основные сведения %%%
\newcommand{\thesisAuthorLastName}{\fixme{Фамилия}}
\newcommand{\thesisAuthorOtherNames}{\fixme{Имя Отчество}}
\newcommand{\thesisAuthorInitials}{\fixme{И.\,О.}}
\newcommand{\thesisAuthor}              % Диссертация, ФИО автора
{%
  \texorpdfstring{% \texorpdfstring takes two arguments and uses the first for (La)TeX and the second for pdf
    \thesisAuthorLastName~\thesisAuthorOtherNames% так будет отображаться на титульном листе или в тексте, где будет использоваться переменная
  }{%
    \thesisAuthorLastName, \thesisAuthorOtherNames% эта запись для свойств pdf-файла. В таком виде, если pdf будет обработан программами для сбора библиографических сведений, будет правильно представлена фамилия.
  }
}
\newcommand{\thesisAuthorShort}         % Диссертация, ФИО автора инициалами
{\thesisAuthorInitials~\thesisAuthorLastName}
\newcommand{\thesisTitle}               % Диссертация, название
{\fixme{Название ВКР}}
\newcommand{\thesisSpecialtyNumber}     % Специальность, номер
{\fixme{03.03.01.}}
\newcommand{\thesisSpecialtyTitle}      % Специальность, название
{\fixme{Прикладные математика и физика (бакалавриат)}}
\newcommand{\thesisProfile}             % Направленность (профиль) подготовки
{\fixme{Авиационные технологии}}
\newcommand{\thesisCity}                % Диссертация, город написания диссертации
{Жуковский}
\newcommand{\thesisYear}                % Диссертация, год написания диссертации
{\fixme{2024}}

\newcommand{\thesisOrganization}        % Организация
{Федеральное государственное автономное образовательное учреждение высшего образования <<Московский физико-технический институт (национальный исследовательский университет)>>}
\newcommand{\thesisOrganizationShort}   % Диссертация, краткое название организации
{МФТИ}
\newcommand{\thesisDepartment}          % Название факультета
{Физтех-школа аэрокосмических технологий}
\newcommand{\thesisDepartmentShort}     % Краткое название факультета
{ФАКТ}
\newcommand{\thesisChair}               % Название кафедры
{Кафедра компьютерного моделирования}

\newcommand{\supervisorFio}             % Научный руководитель, ФИО
{\fixme{Фамилия Имя Отчество}}
\newcommand{\supervisorRegalia}         % Научный руководитель, регалии
{\fixme{кандидит технических наук}}
\newcommand{\supervisorFioShort}        % Научный руководитель, ФИО
{\fixme{Фамилия\:И.\,О.}}
\newcommand{\supervisorRegaliaShort}    % Научный руководитель, регалии
{\fixme{к.\,т.\,н.}}

% To avoid conflict with beamer class use \providecommand
\providecommand{\keywords}%             % Ключевые слова для метаданных PDF диссертации и автореферата
{}
